\documentclass[gmd, manuscript]{copernicus}

\graphicspath{{assets/}}
\setlength{\emergencystretch}{3em}

\begin{document}
\sloppy

\title{py2sess v0.2.0: a Python platform for differentiable 2 stream-exact single scattering radiative-transfer benchmarks}

\Author[1][author@email.edu]{Author}{Name}
\Author[1]{Coauthor}{Name}

\affil[1]{Institution, Department, Address}

\runningtitle{Differentiable 2S-ESS benchmarks with py2sess}
\runningauthor{Author et al.}

\received{}
\pubdiscuss{}
\revised{}
\accepted{}
\published{}
\firstpage{1}

\maketitle

\begin{abstract}
Atmospheric retrieval development needs radiative-transfer (RT) solvers whose spectra,
timings, and sensitivities can be tested in the same scientific workflow.  py2sess
v0.2.0 maps the 2 stream-exact single scattering (2S-ESS) method into NumPy and
PyTorch without changing the underlying approximation.  It reproduces exported Fortran
2S-ESS thermal-emission and solar-scattering benchmark spectra with maximum relative
differences below 0.005 \%.  Prepared-input timing experiments then isolate the RT
kernel from opacity generation.  Forward runtime is controlled by wavelength count,
vertical depth, source-regime complexity, and backend throughput.  PyTorch
vector--Jacobian products (VJPs) make local layer and surface sensitivities available
from the tested forward code, but they are not free: on an NVIDIA A100, 50 000
wavelengths and 114 layers require 0.26--0.77 s for wavelength-local surface
sensitivities and 0.52--2.25 s for layer-resolved optical-property sensitivities.
Activating one layer or all layers changes these VJP timings little, indicating that
graph retention and reverse traversal dominate over a loop over state variables.  A
small OCO-3 reflected-solar example demonstrates the same Python workflow on measured
spectra under constrained aerosol and spectroscopy assumptions.  The result is a
validated platform for prepared-input 2S-ESS benchmarking, retrieval prototyping, and
differentiable sensitivity studies.
\end{abstract}

\introduction

Radiative-transfer (RT) models are the forward operators in atmospheric spectroscopy,
remote sensing, data assimilation, and optimal-estimation retrievals.  Retrieval
algorithms need both spectra and derivatives with respect to atmospheric and surface
state variables; this role of the forward-model Jacobian is central to the inverse
problem formulation of Rodgers (2000) and to linearized RT solvers such as LIDORT and
VLIDORT (Spurr, 2006; Spurr and Natraj, 2011; Spurr et al., 2022).  Fast forward RT is
therefore necessary but not sufficient.  For retrieval development and sensitivity
analysis, the practical bottleneck is often the ability to obtain reliable derivatives
for a changing state-vector definition without rewriting and revalidating a new
hand-linearized model each time.

This issue has become more important as retrieval prototyping, uncertainty analysis,
observing-system simulation, and machine-learning emulator development have moved into
Python array workflows.  NumPy provides a common CPU array interface for scientific
computing (Harris et al., 2020), while PyTorch provides accelerator execution and
automatic differentiation in the same programming model used for machine learning
(Baydin et al., 2018; Paszke et al., 2019).  In RT, differentiable implementations are
increasingly used to connect forward models with inverse modeling and learning-based
workflows, including recent differentiable Monte Carlo and three-dimensional RT studies
(Salesin et al., 2024; Wang et al., 2024; Cao et al., 2025).  A one-dimensional,
hyperspectral RT solver that can be called, timed, and differentiated directly in this
environment is therefore useful even when mature Fortran solvers remain the reference
for production calculations.

The 2 stream-exact single scattering (2S-ESS) method combines a two-stream
multiple-scattering approximation with exact single scattering in a curved atmosphere
(Spurr and Natraj, 2011; Natraj et al., 2023).  The method is designed for
high-throughput spectral radiance calculations: the multiple-scattering field is
approximated efficiently, while the single-scattering term retains accurate phase
function and slant-path geometry.  The Fortran 2S-ESS implementation is optimized and
analytically linearized, and 2S-ESS has been used in greenhouse-gas, aerosol, and
temperature/composition retrieval studies (Natraj et al., 2022, 2023).  The present
paper does not introduce a new RT approximation or claim to supersede that analytic
Fortran linearization.  Instead, it asks a platform question: what prepared-input RT
performance and derivative cost are obtained when the same 2S-ESS algebra is exposed
through a Python implementation with NumPy and PyTorch backends?

This benchmark is complementary to RT acceleration and application studies that reduce
the cost of existing retrieval forward models through principal-component or
lookup-table ideas (Spurr et al., 2013; Bak et al., 2021; Natraj and Spurr, 2025).  The
question here is narrower but different: once prepared optical properties and source
arrays are available, how does a Python-native 2S-ESS kernel scale, and what additional
cost is introduced when local sensitivities are obtained by PyTorch reverse-mode
differentiation?

The answer matters scientifically because state-vector experimentation is a real part
of retrieval research.  A retrieval developer may need to switch sensitivities among
optical thickness, single-scattering albedo, asymmetry parameter, thermal source terms,
surface albedo, or instrument-specific surface parameters.  With a differentiable
backend, these experiments can be performed in the same tested forward model and
checked against finite differences.  This flexibility is also relevant to emulator and
surrogate development, where local RT sensitivities may help diagnose training data or
couple differentiable RT components with learned models.  The central quantity for such
work is not a language-level speedup, but a reproducible cost map for forward RT and
local sensitivity calculations after optical properties are available.

We report three benchmark groups and one real-data workflow demonstration.  First,
exported full-spectrum thermal-emission and solar-scattering benchmark scenes
compare py2sess spectra against the Fortran 2S-ESS reference and measure RT
module/component timing for Fortran, NumPy, PyTorch CPU, and PyTorch CUDA.  Second,
synthetic forward-only cases use deterministic inhomogeneous prepared optical profiles
to isolate scaling with wavelength count and vertical-grid size.  Third, synthetic
PyTorch cases measure vector--Jacobian-product (VJP) timings for local prepared-input
sensitivities.
Finally, a small OCO-3 reflected-solar example tests whether the same Python workflow
can ingest measured spectra after gas optical-depth caches and instrument inputs have
been prepared.  This OCO-3 example is demoted deliberately to a workflow
demonstration: it uses open spectroscopy and constrained aerosol assumptions and is not
an operational aerosol product or an independent validation of the OCO-3 Level-2
retrieval.  The solver regimes are thermal emission and solar-scattering observation;
the O$_2$ A-band benchmark is described as Solar rather than ultraviolet.  GPU
results are reported for NVIDIA Tesla T4 and A100 systems where CUDA was available.
The benchmark study is deliberately limited to the RT kernel: opacity generation, file
input/output, scene parsing, and host-to-device tensor construction are outside the
timed region.

\section{Model and benchmark design}

\subsection{2S-ESS radiative-transfer equations}

py2sess follows the 2S-ESS decomposition described by Natraj et al. (2023).  The
upwelling radiance at the top of the atmosphere is written as the sum of a two-stream
multiple-scattering contribution and an exact first-order contribution,
\begin{equation}
  I_{\rm TOA}^{\uparrow}
  =
  I_{\rm 2S}^{\uparrow}
  +
  I_{\rm FO}^{\uparrow}.
  \label{eq:rt_decomposition}
\end{equation}
The two-stream term solves for the diffuse field on the layer grid, including
regime-specific particular sources.  The FO term is evaluated outside the two-stream
boundary-value problem and added at the same output level.  In thermal-emission cases,
the two-stream calculation is azimuth independent and uses Planck-source particular
terms; the FO component is the exact line-of-sight thermal and surface-emission term.
In solar-scattering cases, the solver evaluates solar slant-path transmission,
direct-beam particular sources, the retained azimuthal Fourier components, and exact
first-order solar scattering.

The prepared RT inputs are layer optical thickness, single-scattering albedo,
asymmetry parameter, source terms, and surface properties.  When optical properties are
constructed from component quantities, py2sess uses
\begin{align}
  \tau_n &= \tau_{n,{\rm abs}} + \tau_{n,{\rm sca}}, \\
  \omega_n &= \frac{\tau_{n,{\rm sca}}}{\tau_n}, \\
  g_n &=
  \frac{\sum_a g_{a,n}\tau_{a,n,{\rm sca}}}
       {\tau_{n,{\rm sca}}},
  \label{eq:optical_inputs}
\end{align}
where Rayleigh scattering contributes zero to $g_n$ and aerosol components contribute
through their scattering optical depths and moments.  For a purely absorbing layer
($\tau_{n,{\rm sca}}=0$), $\omega_n=0$ and $g_n$ is set to zero by convention because
the asymmetry parameter then has no radiative effect.  The solar FO source uses the
corresponding mixture of Rayleigh and aerosol phase functions at the scattering angle.

The strongly forward-scattering part is removed before the two-stream solve by the
delta-M transformation,
\begin{equation}
\begin{aligned}
\widetilde{\tau}_n &= (1-\omega_n f_n)\tau_n,\\
\widetilde{\omega}_n &= \frac{(1-f_n)\omega_n}{1-\omega_n f_n},\\
\widetilde{g}_n &= \frac{g_n-f_n}{1-f_n},
\end{aligned}
\label{eq:deltam}
\end{equation}
where $f_n$ is the truncation factor.  The default Henyey--Greenstein-like choice is
$f_n=g_n^2$, while mixed Rayleigh--aerosol inputs may provide $f_n$ from phase
moments.  In the derivative benchmarks, $f_n$ is treated as a prepared RT input; the
reported sensitivities do not include derivatives through phase-table construction or
through the default $f_n=g_n^2$ rule.  Numerical endpoint regularization is applied to
avoid exactly singular two-stream limits.

For a homogeneous layer and Fourier component $m$, the two-stream homogeneous
equations can be written after the standard sum--difference transformation as
\begin{equation}
  \frac{{\rm d}}{{\rm d}\tau}
  \begin{bmatrix}
    y_{m,1}\\
    y_{m,2}
  \end{bmatrix}
  =
  \begin{bmatrix}
    0 & s_m\\
    d_m & 0
  \end{bmatrix}
  \begin{bmatrix}
    y_{m,1}\\
    y_{m,2}
  \end{bmatrix}
  +
  \mathbf{q}_m(\tau),
  \label{eq:twostream_system}
\end{equation}
where $\tau$ is optical depth measured downward within the layer and
$(y_{m,1},y_{m,2})$ are the sum--difference variables equivalent to upwelling and
downwelling quadrature intensities.  The two-stream ordinate $\bar{\mu}$ is a closure
choice.  Classical two-stream formulations include quadrature choices such as
$\bar{\mu}=1/\sqrt{3}$ and hemispheric or mean-value closures with different angular
moments (Meador and Weaver, 1980; Toon et al., 1989; Spurr and Natraj, 2011).  The
benchmarks use the ordinates in the exported reference settings:
$\bar{\mu}=0.5$ for the thermal-emission case and $\bar{\mu}=1/\sqrt{3}$ for the
solar-scattering case.  With
$c_n=3\widetilde{\omega}_n\widetilde{g}_n$,
$p_0^2=\bar{\mu}^2$, and $p_1^2=(1-\bar{\mu}^2)/2$, the retained Fourier components
use
\begin{equation}
\begin{alignedat}{2}
s_{0,n} &= \frac{\widetilde{\omega}_n-1}{\bar{\mu}},
&\quad d_{0,n} &= \frac{p_0^2c_n-1}{\bar{\mu}},\\
s_{1,n} &= \frac{p_1^2c_n-1}{\bar{\mu}},
&\quad d_{1,n} &= -\frac{1}{\bar{\mu}} .
\end{alignedat}
\label{eq:twostream_sd}
\end{equation}
Thermal-emission calculations use only $m=0$; solar-scattering calculations retain the
$m=0$ and $m=1$ terms used by the reference benchmark.

The homogeneous eigenvalue and transformed up/down eigenvectors are
\begin{equation}
\begin{aligned}
  k_{m,n} &= \sqrt{s_{m,n}d_{m,n}},\\
  \chi_{m,n} &= -\frac{s_{m,n}}{k_{m,n}},\\
  \mathbf{X}^{+}_{m,n}
    &=\frac{1}{2}
      \begin{bmatrix}1+\chi_{m,n}\\1-\chi_{m,n}\end{bmatrix},\\
  \mathbf{X}^{-}_{m,n}
    &=\frac{1}{2}
      \begin{bmatrix}1-\chi_{m,n}\\1+\chi_{m,n}\end{bmatrix}.
\end{aligned}
\label{eq:twostream_eigen}
\end{equation}
The layer solution in quadrature-intensity space is then
\begin{equation}
  \begin{split}
  \mathbf{I}_{m,n}^{\rm 2S}(\tau)
  &=
  L_{m,n}\mathbf{X}_{m,n}^{+}e^{-k_{m,n}\tau}\\
  &\quad
  +
  M_{m,n}\mathbf{X}_{m,n}^{-}e^{-k_{m,n}(\widetilde{\tau}_n-\tau)}
  +
  \mathbf{W}_{m,n}(\tau),
  \end{split}
  \label{eq:twostream_solution}
\end{equation}
where $L_{m,n}$ and $M_{m,n}$ are layer constants and $\mathbf{W}_{m,n}$ is the
particular solution.

For the solar two-stream particular solution, the py2sess solar batch path used in
these benchmarks evaluates the direct-beam source with the Green-function beam
multiplier treatment described for 2S-ESS by Natraj et al. (2023).  Within a layer the
source has the form
\begin{equation}
  \mathbf{q}_{m,n}^{\odot}(\tau)
  =
  \mathbf{Q}_{m,n}^{\odot}\exp(-\eta_{n}\tau),
  \qquad
  \gamma_{\pm,m,n}=\eta_n\pm k_{m,n},
  \label{eq:solar_green}
\end{equation}
where $\eta_n$ is the layer-average solar secant from the slant-path geometry.  Taylor
limits are used when the denominators involving $\gamma_{\pm,m,n}$ are small.
$\mathbf{Q}_{m,n}^{\odot}$ contains the Fourier normalization, beam flux, phase-moment
terms, and layer optical-property factors for the selected source term.  For the
two-stream thermal source, Planck's law is evaluated at the relevant wavelength or
wavenumber and represented linearly in each layer,
\begin{equation}
  B_n(\tau)=B_{n,0}+B_{n,1}\tau,\qquad
  \mathbf{q}_{n}^{\rm th}(\tau)
  =
  (1-\widetilde{\omega}_n)B_n(\tau)
  \begin{bmatrix}1\\1\end{bmatrix}.
  \label{eq:thermal_source}
\end{equation}
The two-stream thermal source in Eq. (\ref{eq:thermal_source}) uses the delta-M-scaled
single-scattering albedo.  The thermal particular solution is evaluated through
analytic source multipliers at layer boundaries and at the user angle.

The boundary-value problem (BVP) for the layer constants is
\begin{equation}
  \begin{aligned}
  \mathbf{A}_m\mathbf{c}_m &= \mathbf{r}_m,\\
  \mathbf{c}_m &=
  (L_{m,1},M_{m,1},\ldots,L_{m,N},M_{m,N})^T .
  \end{aligned}
  \label{eq:bvp}
\end{equation}
Rows enforce no incoming diffuse radiation at the top boundary, continuity across
internal boundaries, and the lower-boundary surface reflection or emission condition.
The mathematical system is banded and block structured because only adjacent layers are
coupled.  With the constants ordered as
$(L_1,M_1,L_2,M_2,\ldots,L_N,M_N)$, each continuity row touches only the two
neighbouring layers, so the regular system has five non-zero diagonals.  The reported
double-precision batch calculations use this pentadiagonal path for both NumPy and
PyTorch.  For each independent wavelength row, py2sess stores the lower-second,
lower-first, main, upper-first, and upper-second coefficient bands, eliminates the two
lower bands in a Thomas-like forward sweep, and recovers the layer constants by
back substitution.  This is an $O(L)$ banded solve per row rather than a dense
$O(L^3)$ solve on the $2L\times2L$ BVP matrix.  For PyTorch VJPs, the BVP solve is
wrapped by a custom autograd rule whose reverse pass solves the transposed
pentadiagonal adjoint system and differentiates the prepared bands and right-hand side.
Dense or block-tridiagonal repair is available for non-finite exceptional rows, but
dense all-row BVP solves are not the primary path in the reported benchmarks.

After the quadrature-angle solution is known, the intensity is post-processed to the
requested user angle $\mu$ through a layer recurrence.  Layers are indexed downward;
$\tau=0$ is the top of layer $n$ and $\tau=\widetilde{\tau}_n$ is its bottom.  For an
upwelling ray, the recurrence marches upward: the bottom-boundary user-angle intensity
is attenuated across the layer and the layer-integrated source contribution is added at
the layer top,
\begin{equation}
  I_{m,n}^{\uparrow}(0,\mu)
  =
  e^{-\widetilde{\tau}_n/\mu}
  I_{m,n}^{\uparrow}(\widetilde{\tau}_n,\mu)
  +
  \mathcal{S}_{m,n}^{\uparrow}(\mu),
  \label{eq:user_post}
\end{equation}
where $\mathcal{S}_{m,n}^{\uparrow}$ contains the analytically integrated homogeneous
and particular source contributions.  The first-order solar term follows the scalar
single-scattering equation in the delta-M-scaled optical-depth coordinate $x$ along the
line of sight,
\begin{equation}
  \begin{aligned}
  \frac{{\rm d}I_n^{\uparrow}(x,\Omega)}{{\rm d}x}
  &=
  -I_n^{\uparrow}(x,\Omega)+S_n^{(1)}(x,\Omega),\\
  S_n^{(1)}
  &=
  \frac{\omega_n P_n(\Omega,\Omega_0)F_0}
       {4\pi(1-\omega_n f_n)}
  \exp[-\widetilde{\tau}_{\odot,n}(x)] .
  \end{aligned}
  \label{eq:ss_rte}
\end{equation}
Here $\widetilde{\tau}_{\odot,n}$ is the solar slant optical depth to the scattering
point, $P_n$ is the mixed Rayleigh--aerosol phase function before the delta-M source
factor, and $F_0$ is the incident solar flux.  Optical depths in the line-of-sight and
solar paths use $\widetilde{\tau}_n$, while the exact-scatter source carries the
corresponding $(1-\omega_n f_n)^{-1}$ factor.  The thermal FO term is the analogous
line-of-sight integration of thermal source and surface emission.  The validated
thermal benchmark follows the exported Fortran reference convention: the FO optical
depth is delta-M scaled, while the local Planck absorption factor remains
$(1-\omega_n)$.  A diagnostic 1000-wavelength, 50-layer synthetic case gives a maximum
radiance change of 0.016 \% when a source-side thermal delta-M multiplier is added, so
this implementation detail is explicit but not a dominant term in the benchmark
scenes.

The implementation mapping is as follows.  Optical-property builders construct
$\tau$, $\omega$, $g$, delta-M factors, and FO solar phase terms from component optical
depths.  Backend-independent delta-M utilities implement Eq. (\ref{eq:deltam}) for
array and tensor inputs.  The thermal and solar batch solvers implement the
homogeneous eigensystem, analytic source multipliers, BVP solve, and user-angle
post-processing in Eqs. (\ref{eq:twostream_system})--(\ref{eq:user_post}).  FO modules
evaluate Eq. (\ref{eq:ss_rte}) and its thermal analogue outside the multiple-scattering
BVP.  The same algebra is dispatched to NumPy arrays or PyTorch tensors.  py2sess adds
vectorized wavelength batching, CUDA-resident tensor execution, and differentiable
PyTorch paths; it does not add custom CUDA kernels or a separate hand-written adjoint
solver in the benchmarks reported here.

Batching follows the public array convention that all leading dimensions before the
final layer axis are independent RT rows.  For a full-spectrum one-column calculation,
the prepared optical-property arrays have shape $(N_\lambda,L)$ and are interpreted as
$N_\lambda$ independent wavelength rows sharing the same layer count and geometry.  For
multi-column or multi-geometry calls, py2sess preserves the leading batch dimensions
and appends a geometry axis only at the output.  Internally the batch solvers flatten
the leading dimensions to a row batch, evaluate layer recurrences in the vertical
dimension, and vectorize the algebra over the row batch.  On CUDA, the prepared tensors
are placed on the device before timing so the measured region covers the RT algebra and
autograd graph traversal, not host-to-device construction.

\subsection{Timing scope}

All timing values in this paper are wall-clock RT times after prepared RT arrays are
available.  File input/output, scene parsing, opacity-table generation, optical-property
construction, and tensor construction are excluded.  For CUDA runs, tensors are already
resident on the target device before the timed region, and CUDA work is synchronized
around each measured repeat.  Reported values use double precision.  Unless otherwise
specified, each timing value is the mean of five measured repeats.  The synthetic
benchmark runs include untimed warmup evaluations before those repeats; the
full-spectrum benchmark rows use the five measured repeats reported by the benchmark
drivers.  The full-spectrum figures show the measured range across repeats; synthetic
figures show the standard deviation.

The local CPU rows were measured on an Apple M2 Pro system.  NumPy was run with
single-threaded BLAS/OpenMP settings.  CUDA rows were measured on NVIDIA Tesla T4 and
A100 GPUs; CPU timings from those accelerator environments are not used in the figures.
Fortran rows are measured from the optimized external 2S-ESS benchmark executable on
the local CPU.  Consequently, the Fortran and Python rows are empirical benchmark-task
timings with different driver scopes, not processor-normalized compiler or language
speedup measurements.  Apple MPS is not included because the benchmark comparison was
defined for NumPy CPU, PyTorch CPU, and PyTorch CUDA.

\subsection{Scope relative to analytic RT Jacobians}

The benchmark is not intended to replace mature hand-linearized RT codes.  Analytic
linearization and tangent-linear/adjoint modules remain the right target for a fixed
operational retrieval state vector because the derivative algebra can be specialized,
validated, and optimized for that product.  Examples include CRTM water-vapor
Jacobian work for numerical weather prediction (Chen et al., 2010),
tangent-linear and adjoint polarized-RT development for polarimetric retrievals
(Ding and Yang, 2023), and the analytic
linearization in the Fortran 2S-ESS reference model (Natraj et al., 2023).  The py2sess
question is narrower: after the 2S-ESS forward algebra has been mapped into Python,
how expensive are local prepared-input derivatives when the differentiated variable is
changed without deriving a new hand-coded Jacobian?  This makes the PyTorch backend a
development and sensitivity-analysis path, not a claim that automatic differentiation
is faster than a tuned analytic Jacobian for a settled production retrieval.

\subsection{Full-spectrum benchmark cases}

The full-spectrum benchmark uses two exported 2S-ESS reference scenes: a
200 000-point thermal-emission spectrum and a 280 000-point solar-scattering spectrum,
each with 114 RT layers.  py2sess reads the prepared RT input bundle derived from these
scenes and computes the Python spectrum with the NumPy backend.  The reference opacity
inputs have HITRAN2020 provenance (Gordon et al., 2022), but opacity generation is
treated as upstream preprocessing and is outside the timed RT region.

\begin{table}[t]
\caption{Full-spectrum benchmark settings.}
\label{tab:full_spectrum_settings}
\centering
\begin{tabular}{lll}
\hline
Quantity & Thermal-emission case & Solar-scattering case \\
\hline
Paper label & Thermal & Solar \\
Spectral interval & 14.29--20.00 $\mu$m & 0.753--0.795 $\mu$m \\
Spectral points & 200 000 & 280 000 \\
RT layers & 114 & 114 \\
Primary gases & H$_2$O, CO$_2$, N$_2$O, O$_3$, HNO$_3$, SO$_2$, NH$_3$ & H$_2$O, O$_2$ \\
Geometry & VZA $=49.51^\circ$ & SZA $=47.70^\circ$, VZA $=49.51^\circ$, RAA $=275.75^\circ$ \\
Surface & emissivity 0.9853, albedo 0.0147, $T_s=300.35$ K &
spectral albedo 0.462--0.492 \\
Two-stream ordinate & $\bar{\mu}=0.5$ & $\bar{\mu}=1/\sqrt{3}$ \\
\hline
\end{tabular}
\end{table}

The spectrum agreement metric is
\begin{equation}
  \epsilon_{\rm spec}
  =
  \max_i
  \left[
    100
    \frac{\left|I_{{\rm py},i}-I_{{\rm F},i}\right|}
         {\max\left(\left|I_{{\rm F},i}\right|,10^{-15}\right)}
  \right],
  \label{eq:spectrum_error}
\end{equation}
where $I_{{\rm F},i}$ is the Fortran reference radiance at spectral point $i$.
The small denominator floor prevents isolated near-zero radiances from dominating the
reported maximum relative error; it is many orders of magnitude below the plotted
radiances in the benchmark spectra.

\subsection{Synthetic forward benchmarks}

Synthetic cases isolate RT scaling from external opacity generation while avoiding a
featureless constant slab.  The solar-scattering case contains a smooth absorbing
background, a low-altitude aerosol layer, and an elevated cloud-like scattering layer
near 0.65 $\mu$m.  The thermal-emission case contains smooth gas-like absorption,
a weak low-altitude scattering layer, surface emissivity 0.98, surface albedo 0.02, and
a deterministic US-standard-like temperature profile for thermal source arrays.  These
profiles are generated analytically by the benchmark script, so no opacity files or
spectroscopic cross sections are used in the timed region.  The delta-M truncation
factor is fixed to $f=0$ in the reported VJP cases to keep the sensitivities local to
the prepared RT variables.

Two one-dimensional sweeps are used.  The spectral sweep varies the number of
wavelengths at fixed 114 RT layers.  The vertical sweep varies the number of RT layers
at fixed 50 000 wavelengths.  This design separates the two dominant dimensions without
requiring a full Cartesian benchmark grid.  For fixed optical properties, wavelengths
are independent while the vertical solution is layer ordered.  We interpret the
measured forward timing with the empirical model
\begin{equation}
  T_{\rm F}^{(b,c)}(N_\lambda,L)
  \simeq
  T_0^{(b)}
  + \kappa_{\rm F}^{(b,c)} N_\lambda C_{\rm F}^{(c)}(L),
  \label{eq:forward_scaling}
\end{equation}
where $N_\lambda$ is the number of wavelengths, $L$ is the number of RT layers,
$T_0^{(b)}$ is backend dispatch and fixed overhead, $C_{\rm F}^{(c)}(L)$ is the
vertical work per wavelength for regime $c$, and $\kappa_{\rm F}^{(b,c)}$ absorbs
backend throughput, vectorization, and memory traffic.  Equation
(\ref{eq:forward_scaling}) is used only to define the dimensions being swept; we do not
fit or interpret it as a formal scaling law.

\subsection{Synthetic VJP sensitivity benchmarks}

The PyTorch sensitivity benchmarks time a VJP: a forward RT evaluation followed by
reverse-mode differentiation of the summed total-radiance objective.  The calculation
records gradients with respect to selected local RT variables but does not materialize
a dense cross-wavelength Jacobian.  The design follows the same validation logic used
in tangent-linear and adjoint RT studies such as Ding and Yang (2023): finite
differences are used for compact correctness checks, while the timed derivative cases
target variables that are scientifically meaningful in retrievals.  In py2sess these
targets are layer optical thickness, single-scattering albedo, asymmetry factor
as a scalar phase-shape proxy, wavelength-local surface albedo, and thermal surface
emissivity.  For the solar omega and $g$ VJPs, the benchmark lets the differentiable
Henyey--Greenstein first-order scatter source respond to the active variable rather
than holding a precomputed FO scatter array fixed.  Delta-M truncation factors remain
fixed, so the sensitivities are still local derivatives with respect to prepared RT
variables rather than end-to-end derivatives through optical-property preprocessing.

Direct py2sess public-API cases were checked against centered finite differences
using the same summed objective and fixed prepared inputs (Table
\ref{tab:gradient_validation}).  The small cases perturb every listed column, and the
medium cases perturb sampled tau, omega, and $g$ columns in 1000-wavelength,
50-layer calculations.  These are prepared-input derivative checks, not a
benchmark-scale dense Jacobian validation.  Relative error is computed as
$|J_{\rm AD}-J_{\rm FD}|/\max(|J_{\rm FD}|,10^{-12})$, where $J_{\rm AD}$ is the
PyTorch automatic-differentiation derivative and $J_{\rm FD}$ is the finite-difference
derivative with perturbation $h=10^{-6}$.  All sampled medium-scale columns have
relative error below $1.3\times10^{-3}$ and no columns above the $10^{-2}$ diagnostic
threshold.  The timing benchmark separately requires finite gradient norms and
checksums for each measured sensitivity row.

The repository also contains compact Fortran-derived Jacobian reference fixtures for
surface-emissivity, surface-temperature, and surface-albedo derivatives.  These
fixtures are used for derivative validation of selected thermal and solar cases, but
they do not contain Fortran analytic-Jacobian runtimes and are therefore not used as
matched timing comparisons.

\begin{table}[t]
  \caption{PyTorch gradient checks against centered finite differences for
  prepared-input RT.}
  \label{tab:gradient_validation}
  \centering
  \scriptsize
  \begin{tabular}{llllllll}
  \hline
  Regime & Case size & Variables & Cols. & Max. $|J_{\rm FD}|$ &
  Max. abs. err. & Mean abs. err. & Max. rel. err. \\
  \hline
  Solar & 2 $\lambda$ $\times$ 3 $L$ & $\tau,\omega,g$ & 18 &
  $2.21\times10^{-2}$ & $3.16\times10^{-12}$ & $9.29\times10^{-13}$ & $1.27\times10^{-8}$ \\
  Solar & 2 $\lambda$ $\times$ 3 $L$ & albedo, flux & 4 &
  $2.45\times10^{-1}$ & $4.86\times10^{-12}$ & $1.81\times10^{-12}$ & $1.03\times10^{-10}$ \\
  Thermal & 2 $\lambda$ $\times$ 3 $L$ & $\tau,\omega,g$ & 18 &
  $2.59\times10^{-1}$ & $3.01\times10^{-10}$ & $1.08\times10^{-10}$ & $2.01\times10^{-8}$ \\
  Thermal & 2 $\lambda$ $\times$ 3 $L$ & Planck terms & 10 &
  $5.05\times10^{-1}$ & $1.67\times10^{-10}$ & $9.65\times10^{-11}$ & $2.97\times10^{-9}$ \\
  Thermal & 2 $\lambda$ $\times$ 3 $L$ & emissivity, albedo & 4 &
  $7.22\times10^{-1}$ & $2.04\times10^{-10}$ & $1.34\times10^{-10}$ & $5.92\times10^{-10}$ \\
  Solar & 1000 $\lambda$ $\times$ 50 $L$ & sampled $\tau,\omega,g$ & 12 &
  $4.83\times10^{-3}$ & $7.04\times10^{-10}$ & $3.04\times10^{-10}$ & $1.29\times10^{-3}$ \\
  Thermal & 1000 $\lambda$ $\times$ 50 $L$ & sampled $\tau,\omega,g$ & 12 &
  $9.15\times10^{-5}$ & $2.45\times10^{-11}$ & $1.03\times10^{-11}$ & $1.25\times10^{-3}$ \\
  \hline
  \end{tabular}
\end{table}

Holding the prepared delta-M factor fixed is also a derivative convention.  In a
200-wavelength, 50-layer diagnostic with $f=g^2$, differentiating through $f(g)$ would
change the Solar $g$-gradient by 1.2 \% in L2 norm and at most 1.7 \% pointwise.
The corresponding thermal $g$-gradient is much smaller in absolute scale; the omitted
chain-rule term is $8.9\times10^{-6}$ in L2 norm.  The benchmark sensitivities should
therefore be read as local derivatives with respect to the prepared RT arrays, not as
end-to-end derivatives through phase-moment or delta-M preprocessing.

The main sensitivity scaling figure uses tau VJPs.  Three sweeps are measured:
wavelength count at fixed 114 layers, layer count at fixed 50 000 wavelengths, and
number of active tau layers at fixed 50 000 wavelengths and 114 layers.  Omega,
$g$, wavelength-local surface-albedo, and thermal surface-emissivity derivatives are
reserved for representative overhead comparisons at the 50 000-wavelength, 114-layer
setting.  The surface cases activate one variable per wavelength, not one scalar shared
by the full spectrum.

We interpret VJP timing with the operational decomposition
\begin{equation}
  T_{{\rm J},q}
  =
  T_{\rm F}
  + T_{{\rm store},q}
  + T_{{\rm rev},q}
  + T_{{\rm leaf},q}(N_g),
  \label{eq:vjp_decomposition}
\end{equation}
where $T_{{\rm J},q}$ is the VJP time for differentiated parameter $q$, $T_{\rm F}$ is
the matching forward RT time, $T_{{\rm store},q}$ is the cost of retaining forward
intermediates needed by derivative rules, $T_{{\rm rev},q}$ is reverse graph traversal,
$T_{{\rm leaf},q}$ is gradient accumulation into active leaf tensors, and $N_g$ is the
number of active local gradient variables.  Equation (\ref{eq:vjp_decomposition}) is a
bookkeeping model for the measured timings; it is not an analytic adjoint cost theorem.

\subsection{OCO-3 aerosol retrieval showcase}

The real-data demonstration uses OCO-3 reflected-solar spectra in the O$_2$ A band
near 0.76 $\mu$m and the weak CO$_2$ band near 1.61 $\mu$m (Eldering et al., 2019).
The strong CO$_2$ band near 2.06 $\mu$m is also built and examined diagnostically, but
it is not used in the reported retrieval because the open spectroscopy used here leaves
larger residual structure in that band.  These OCO bands are used by the ACOS/OCO
full-physics retrieval because the O$_2$ A band constrains photon path length and
scattering, while the CO$_2$ bands provide complementary absorption and surface
information (O'Dell et al., 2018; Bell et al., 2023).  We use OCO-3 Level-1B
radiances and geometry, Level-2 atmospheric profiles and aerosol component labels, and
open HITRAN2020 spectroscopy to build cached gas optical-depth inputs.  Official OCO
ABSCO tables are not used.  The retrieval is therefore a py2sess capability test with
open spectroscopy, not a replacement for the operational ACOS forward model.  During
optimization, gas optical depths above 80 are clipped only in saturated line cores; at
that opacity the transmission is already negligible, and the clipping avoids
undefined reverse-mode products in the O$_2$ A-band calculation.

For the reported pilot case, nine land soundings were selected from one OCO-3 granule
on 24 June 2022 near 40$^\circ$ N, 60$^\circ$ E.  The screen required finite L1B/L2
geometry and radiances, valid L1B--L2 sounding matching, OCO-3 quality flag zero,
land fraction at least 50 \%, O$_2$ A-band Level-2 relative RMS below 0.35 \%, and a
local high-AOD patch within 50 km.  The retrieval state is intentionally small: one
shared aerosol optical depth at 0.755 $\mu$m, band-dependent continuum terms, and
small wavelength shift/stretch nuisance terms.
The OCO-3 Level-2 AOD is not used as a prior.  It is used only as a same-sounding
operational reference.  External aerosol context is supplied by MAIAC MCD19A2 AOD and
MERRA-2 aerosol diagnostics (Gelaro et al., 2017; Lyapustin et al., 2018).  AERONET
Version-3 AOD was checked, but no Level-2 station collocation within the local patch
was available for this pilot scene (Giles et al., 2019).

The pilot scene is dust dominated in the collocated MERRA-2 three-dimensional aerosol
profile.  We therefore use a dust-like effective aerosol model with Henyey--Greenstein
asymmetry factor $g=0.85$, single-scattering albedo $\omega=0.94$, and Angstrom
exponent $\alpha=0.20$.  These fixed values define an effective optical model for the
showcase, not a microphysical aerosol retrieval.  The vertical distribution is
prescribed from the MERRA-2 layer aerosol mass profile rather than retrieved.  If
$m_n$ is the non-negative
MERRA-2 aerosol mass mapped to RT layer $n$, the normalized aerosol layer weight is
\begin{equation}
  w_n
  =
  \frac{m_n}{\sum_j m_j},
  \label{eq:oco_aerosol_profile}
\end{equation}
which fixes the aerosol profile shape while allowing the column AOD to be inferred
from the OCO-3 spectra.  The aerosol optical thickness in band $b$ is then
\begin{equation}
  \tau_{{\rm aer},b,n}(\lambda)
  =
  {\rm AOD}_{755}
  \left(\frac{\lambda}{0.755~\mu{\rm m}}\right)^{-\alpha}
  w_n ,
  \label{eq:oco_aerosol_tau}
\end{equation}
with fixed Angstrom exponent $\alpha$.  AOD at 0.55 $\mu$m is reported for comparison
with MAIAC and MERRA-2.  Since the MERRA-2 profile shape is used as a prior, the
MERRA-2 column AOD comparison is interpreted as a consistency check rather than a
fully independent validation target.

The retrieval minimizes a weighted least-squares cost over the retained channels,
\begin{equation}
  \Phi(\mathbf{x})
  =
  \frac{1}{2}
  \left[\mathbf{y}-\mathbf{F}(\mathbf{x})\right]^T
  \mathbf{S}_{\epsilon}^{-1}
  \left[\mathbf{y}-\mathbf{F}(\mathbf{x})\right]
  +
  \frac{1}{2}
  \left[\mathbf{x}-\mathbf{x}_{a}\right]^T
  \mathbf{S}_{a}^{-1}
  \left[\mathbf{x}-\mathbf{x}_{a}\right],
  \label{eq:oco_cost}
\end{equation}
where $\mathbf{y}$ is the measured spectrum, $\mathbf{F}$ is the py2sess forward model
with prepared gas and aerosol inputs, and $\mathbf{x}$ contains log-AOD, continuum, and
wavelength nuisance parameters.  The nuisance terms absorb continuum and dispersion
errors that are not part of the 2S-ESS RT kernel; this improves fit stability but also
creates degeneracy with aerosol amount.  For this reason the OCO-3 result is presented
as a constrained workflow demonstration rather than as an aerosol validation product.

\section{Results}

\subsection{Full-spectrum agreement and module/component timing}

py2sess agrees with the Fortran full-spectrum benchmark spectra while exposing a
different timing profile.  The Python and Fortran spectra overlap at plotting
precision for both benchmark scenes (Fig. \ref{fig:full_spectrum_spectra}).  The
maximum relative differences from Eq. (\ref{eq:spectrum_error}) are
$5.0\times10^{-4}\%$ for the thermal case and $4.7\times10^{-3}\%$ for the
Solar case.  These accuracy values are used only to verify that the benchmark
timing is being measured for the intended RT calculation.

\begin{figure}[t]
  \centering
  \includegraphics[width=14cm]{full_spectrum_spectrum_comparison.png}
  \caption{Full-spectrum agreement between py2sess and the Fortran 2S-ESS reference.
  Panels (a) and (c) show the thermal-emission and solar-scattering spectra,
  respectively; panels (b) and (d) show the corresponding Python-minus-Fortran
  relative differences.  Curves are decimated for display; agreement metrics are
  computed over the full spectral grids.}
  \label{fig:full_spectrum_spectra}
\end{figure}

On the Apple M2 Pro, the NumPy thermal component-timing median was 1.16 s, with a repeat
range of 1.09--2.75 s; the mean shown in Fig. \ref{fig:full_spectrum_runtime} is
1.46 s because one repeat was slower than the others.  The NumPy Solar median
was 3.16 s, with a range of 3.13--3.41 s and a mean of 3.21 s.  The optimized Fortran
driver-reported module means were 2.50 s for the thermal case and 27.1 s for the
Solar case on the same machine.  PyTorch CPU required mean times of 3.47 s and
8.54 s.  CUDA
component timings were shorter for these large spectral grids: the A100 required
0.114 s for the thermal case and 1.73 s for the Solar case, while the Tesla T4
required 0.690 s and 2.31 s.

These numbers should be interpreted through the timing scope in Sect. 2.2.  The Python
bars are FO+2S component timings on prepared arrays, while the Fortran bars are
driver-reported RT module timings from the external benchmark executable.  The main
result is therefore not a language-normalized speedup.  Rather, the same validated RT
calculation can be run in three useful py2sess modes: practical NumPy CPU forward
execution, differentiable PyTorch execution, and high-throughput CUDA execution.

\begin{figure}[t]
  \centering
  \includegraphics[width=14cm]{full_spectrum_rt_runtime.png}
  \caption{Full-spectrum RT timing.  Python bars show FO+2S component timing on
  prepared arrays; Fortran bars show the external driver-reported RT module time.
  Values are means of five runs, and error bars show the measured range.  Because the
  timer scopes differ, the Fortran and Python bars should not be interpreted as
  language-normalized speedup ratios.}
  \label{fig:full_spectrum_runtime}
\end{figure}

The Solar case is slower than the thermal case in both the full-spectrum and synthetic
benchmarks.  This is expected from the algorithmic structure.  The thermal-emission
case is azimuth independent and has thermal source terms.  The solar-scattering case
adds solar slant-path transmission, direct-beam source multipliers, two retained
azimuthal Fourier components, and exact first-order solar scattering.  The extra
geometry and source terms increase the per-wavelength vertical work in Eq.
(\ref{eq:forward_scaling}).

\subsection{Forward RT scaling}

\begin{figure}[t]
  \centering
  \includegraphics[width=14cm]{paper_rt_synthetic_forward_publication.png}
  \caption{Synthetic forward RT scaling.  The spectral-grid sweep varies the number of
  wavelengths at fixed 114 RT layers.  The vertical-grid sweep varies the number of RT
  layers at fixed 50 000 wavelengths.  Error bars show the standard deviation from
  five runs.}
  \label{fig:synthetic_forward}
\end{figure}

Synthetic optical-property cases isolate RT scaling from opacity preparation.  With
114 RT layers fixed, forward runtime grows with the number of wavelengths after fixed
backend costs are amortized (Fig. \ref{fig:synthetic_forward}).  At 300 000
wavelengths, the thermal runtime was 1.76 s for NumPy, 5.79 s for PyTorch CPU,
1.10 s on the Tesla T4, and 0.188 s on the A100.  The corresponding Solar
runtimes were 3.46 s, 9.94 s, 2.42 s, and 1.72 s.  At small spectral grids, CUDA
launch and framework overheads limit the accelerator advantage.  At large spectral grids, the accelerator
backends become throughput limited and the A100 gives the shortest RT time.

The vertical dimension has a stronger effect.  At fixed 50 000 wavelengths, increasing
the grid from 5 to 200 RT layers increased the thermal runtime from 0.016 s to
0.518 s for NumPy, from 0.045 s to 1.72 s for PyTorch CPU, from 0.015 s to
0.322 s on the Tesla T4, and from 0.0087 s to 0.088 s on the A100.  The
Solar case reached 1.30 s, 2.83 s, 1.15 s, and 0.963 s at 200 layers on the
same backends.  These trends show that
$C_{\rm F}^{(c)}(L)$ in Eq. (\ref{eq:forward_scaling}) is not a constant-time vector
operation.  The solver propagates coupled quantities through the vertical column and
evaluates layer source terms, so the measured layer response includes the recurrence,
cache and memory traffic, and backend vectorization effects associated with a deeper
column.

\subsection{VJP sensitivity scaling}

The differentiable timing is strongly associated with the size of the RT computational
graph.  For tau VJP timings at 114 RT layers, the Apple M2 Pro runtime increased from
0.090 s to 2.07 s for the thermal case and from 0.247 s to 6.04 s for the
Solar case as the spectral grid increased from 300 to 10 000 wavelengths
(Fig. \ref{fig:synthetic_jacobian}).  Over the same spectral range, the A100 runtime
changed only from 0.361 s to 0.379 s for the thermal case and from 0.923 s to
0.949 s for the Solar case, indicating that fixed PyTorch and accelerator
overheads remain important at these grid sizes.  The VJP wavelength sweep stops at
10 000 wavelengths because reverse-mode timing retains and traverses graph
intermediates that are absent from the forward-only timing; the forward-only sweep can
therefore be extended farther.

Layer scaling is much larger.  At 50 000 wavelengths, increasing from 5 to 200 RT
layers raised the thermal tau-sensitivity runtime from 0.126 s to 28.0 s on the
Apple M2 Pro, from 0.040 s to 5.87 s on the Tesla T4, and from 0.031 s to 1.15 s on
the A100.  The Solar runtime increased from 0.210 s to 82.1 s, from 0.061 s to
17.7 s, and from 0.063 s to 6.69 s on the same systems.  The vertical recurrence and
source terms therefore dominate the reverse-mode pass.  The same Solar--thermal
separation seen in forward RT persists because the reverse pass must retain and
traverse the larger solar-source graph.

Endpoint log-slope exponents summarize the one-dimensional sweeps.  These are
descriptive slopes between the smallest and largest measured cases, not fitted scaling
laws.  For backend $b$ and case $c$, let $x$ be one swept variable: the number of
wavelengths $N_\lambda$, the number of RT layers $L$, or the number of active local
gradient variables $N_g$.  We report
\begin{equation}
  a_x^{(b,c)}
  =
  \frac{\ln T_{\rm J}^{(b,c)}(x_{\max})-\ln T_{\rm J}^{(b,c)}(x_{\min})}
       {\ln x_{\max}-\ln x_{\min}},
  \qquad x\in\{N_\lambda,L,N_g\}.
  \label{eq:vjp_exponent}
\end{equation}
The tau-sensitivity wavelength exponent is nearly linear on the Apple M2 Pro:
$a_{N_\lambda}=0.89$ for the thermal case and $0.91$ for the Solar case.  It is
much smaller on the GPUs over the sampled 300--10 000 wavelength range
($a_{N_\lambda}=0.01$--0.16), where fixed framework and kernel-dispatch costs remain
visible.  The layer exponent is larger:
$a_L=1.35$--1.62 on the Apple M2 Pro and Tesla T4, and $a_L=0.98$--1.26 on the A100.
In contrast, the active-gradient-state exponent is negligible in this experiment,
spanning only $-0.004 \le a_{N_g} \le 0.008$ across platforms and regimes.

At 50 000 wavelengths and 114 layers, the number of active tau-gradient variables
increases from $5.0\times10^{4}$ to $5.7\times10^{6}$ across the gradient-state sweep,
yet the runtime remains nearly constant.  For the thermal case, the PyTorch CPU,
Tesla T4, and A100 runtimes remain near 11.0 s, 2.27 s, and 0.52 s.  For the
Solar case, they remain near 29 s, 6.05 s, and 1.93 s.  Requesting all local
layer sensitivities is therefore not equivalent to repeating the forward RT solve for
each layer.

\begin{figure}[t]
  \centering
  \includegraphics[width=14cm]{paper_rt_synthetic_jacobian_publication.png}
  \caption{Synthetic tau VJP/local-sensitivity RT scaling.  The left column varies the
  number of wavelengths up to 10 000 at fixed 114 RT layers; this VJP sweep is shorter
  than the forward-only spectral sweep because graph retention is required.  The middle
  column varies the number of RT layers at fixed 50 000 wavelengths.  The right column
  varies the number of active tau-gradient variables at fixed 50 000 wavelengths and
  114 RT layers.  Error bars show the standard deviation from five runs.}
  \label{fig:synthetic_jacobian}
\end{figure}

\subsection{VJP sensitivity overhead}

In Fig. \ref{fig:jacobian_overhead}, overhead means the multiplicative ratio between a
VJP sensitivity timing and the matching forward-only RT timing for the same
50 000-wavelength, 114-layer synthetic case:
\begin{equation}
  R_q
  =
  \frac{T_{{\rm J},q}}{T_{\rm F}}
  \simeq
  1 +
  \frac{T_{{\rm store},q}+T_{{\rm rev},q}+T_{{\rm leaf},q}}
       {T_{\rm F}} .
  \label{eq:vjp_overhead}
\end{equation}
Thus $R_q=3$ means that the differentiated run takes three times as long as the matching
forward run.

A wavelength-local surface-albedo derivative costs 2.3--5.1 times the matching forward
RT calculation.  Layer-resolved tau derivatives cost 10.4--12.8 times forward RT for
the thermal case and 6.7--17.7 times forward RT for the Solar case.
Layer-resolved omega derivatives are more expensive, reaching 13.6--18.7 times
forward RT for the thermal case and 7.7--25.6 times forward RT for the Solar
case.  For both tau and omega, the one-layer and all-layer timings are nearly
identical.  The extra cost is therefore consistent with graph retention and reverse
traversal, not with a linear loop over requested layer variables.

The difference among derivative targets follows the graph in Eq.
(\ref{eq:vjp_decomposition}).  The surface-albedo variable enters mainly through
boundary and source terms.  Tau participates in layer transmittances, homogeneous
multipliers, source integrals, and BVP coupling.  Omega changes the local scattering
operator: scaled single-scattering albedo, eigenvalue and eigenvectors, user-angle
amplitudes, and particular-source coefficients.  In thermal cases it also changes the
absorption--scattering partition of the thermal source.  In these synthetic overhead
cases $f=0$, so omega does not affect scaled optical thickness through the delta-M
factor; its larger overhead comes from the scattering and source graph rather than from
that pathway.

\begin{figure}[t]
  \centering
  \includegraphics[width=14cm]{paper_rt_synthetic_jacobian_overhead_publication.png}
  \caption{VJP sensitivity overhead relative to forward RT.  Ratios are shown for the
  tested 50 000-wavelength, 114-layer synthetic cases.  Surface albedo is
  wavelength-local, with one active variable per wavelength.  Tau and omega are
  evaluated with one active layer or all 114 active layers.  The figure reports VJP
  timings, not dense Jacobian materialization.}
  \label{fig:jacobian_overhead}
\end{figure}

\subsection{OCO-3 aerosol retrieval showcase}

The OCO-3 pilot case tests whether the prepared-input py2sess workflow can be used as
a real retrieval-development environment, not only as a synthetic RT benchmark.  The
case is deliberately constrained: gas optical depths are built from open spectroscopy,
the aerosol vertical shape is prescribed, and only a shared aerosol column amount plus
spectral nuisance terms are adjusted.  This design makes the case useful as a platform
demonstration, but it also makes the limits auditable.

With the MERRA-2 profile weights, $g=0.85$, $\omega=0.94$, and $\alpha=0.20$, the
O$_2$ A-band + weak-CO$_2$ retrieval gives a mean AOD550 of 0.735 for the nine
soundings (Fig. \ref{fig:oco3_aerosol_external}).  The colocated reference means are
0.834 for MAIAC 1 km daily AOD550 and 0.747 for MERRA-2 AOD550.  Relative to MAIAC,
the py2sess retrieval has bias -0.100, mean absolute difference 0.174, and RMSE
0.213.  Relative to MERRA-2, the bias is -0.013, the mean absolute difference is
0.108, the RMSE is 0.139, and the Pearson correlation is 0.857.  The
continuum-normalized spectral residual RMS is 1.12 \%.

The comparison should be interpreted conservatively.  The MERRA-2 column AOD is not a
fully independent validation target because the MERRA-2 vertical profile supplies the
retrieval profile shape.  It is nevertheless a useful consistency check: using the
MERRA-2 profile moves the mean py2sess AOD550 from 0.598 in the earlier Gaussian-profile
run to 0.735 and reduces the MERRA-2 mean absolute difference from 0.155 to 0.108.
The MAIAC field is nearly flat over this small patch, whereas the OCO-3 spectra and
MERRA-2 retain a local gradient, so the MAIAC comparison is best read as regional
magnitude consistency rather than pixel-level spatial validation.

The strong CO$_2$ band is the main unresolved limitation in this pilot case.  Including
it in a shared three-band retrieval increases the residual RMS and pulls the aerosol
column lower, indicating that the open HITRAN-based ``ABSCO-lite'' inputs are not yet
adequate for using that band as a quantitative aerosol constraint.  The reported
retrieval therefore uses O$_2$ A plus weak CO$_2$, and the strong-CO$_2$ result is kept
as a diagnostic rather than as the paper-facing aerosol estimate.  The outcome is a
constrained real-data retrieval showcase: it demonstrates that py2sess can connect
measured OCO-3 spectra, open gas spectroscopy, a simple aerosol state vector, and
external aerosol context in one reproducible workflow.  It should not be described as
an operational aerosol validation.

\begin{figure}[t]
  \centering
  \includegraphics[width=14cm]{oco3_patch_external_comparison.png}
  \caption{OCO-3 aerosol pilot-case comparison for the 24 June 2022 dust patch.  The
  py2sess retrieval uses the O$_2$ A band and weak CO$_2$ band with a MERRA-2
  vertical-profile constraint and no MERRA-2 column-AOD prior.  Panel (a) compares
  retrieved AOD550 with MAIAC, MERRA-2, and the earlier Gaussian-profile py2sess
  baseline.  Panel (b) compares retrieved AOD550 with the external AOD550 references.
  Panel (c) maps the retrieved AOD550 over the local OCO-3 patch.}
  \label{fig:oco3_aerosol_external}
\end{figure}

\section{Discussion}

These benchmarks support a platform contribution, not a new RT approximation.  The
2S-ESS equations in Sect. 2.1 are the published 2S-ESS formulation mapped into a
Python implementation with array and tensor backends.  The full-spectrum benchmarks
show that this mapping reproduces the Fortran reference spectra at the precision needed
for the timing study.  The synthetic benchmarks then show where the cost lies when
opacity generation is held fixed: forward RT is controlled by wavelength count,
vertical depth, source-regime complexity, and backend throughput, while VJP timing is
controlled by the differentiated graph that must be stored and traversed.

The Solar case is generally slower than the thermal case because the solar-scattering regime
contains more physics and more algebra per wavelength.  Solar slant paths, direct-beam
particular sources, retained Fourier terms, and exact first-order scattering all add to
the vertical work.  This distinction is visible even in synthetic cases generated from
the same analytic optical-profile recipe, so it should be treated as a solver-regime
effect rather than as an artifact of a particular opacity table.

The VJP results clarify the statement that automatic differentiation can make
Jacobians easy to obtain.  It does not make them free.  Layer-resolved optical-property
sensitivities cost several to tens of forward RT solves in the tested cases, and a
hand-linearized solver can remain the better production choice for a fixed,
well-established state vector.  This is consistent with the broader RT literature:
operational fast models and retrieval RTMs often invest in tangent-linear, adjoint, or
analytic Jacobian modules because those derivatives are central to data assimilation
and retrieval performance (Chen et al., 2010; Ding and Yang, 2023).  The Fortran 2S-ESS
reference is also analytically linearized (Natraj et al., 2023), so a matched Fortran
Jacobian benchmark would likely remain faster for the exact variables already
implemented there.  The A100 timings here are nevertheless in the same practical
regime for development-size prepared-input tests: at 50 000 wavelengths and 114 layers,
surface-albedo VJPs take less than 1 s and tau or omega VJPs take about
0.5--2.3 s depending on regime.  The advantage of py2sess is therefore not replacing a
tuned analytic Jacobian in production; it is that the same tested RT implementation can
expose sensitivities for new local variables without deriving and maintaining a new
analytic linearization.  Once the differentiated graph is fixed, the number of active
local layer variables has little effect on runtime because the VJP is not a set of
repeated forward perturbations and the benchmark does not materialize a dense
cross-wavelength Jacobian.

This cost structure can support retrieval prototyping and machine-learning-adjacent
workflows.  During retrieval development, the differentiated variable can be changed
from optical thickness to scattering albedo, asymmetry, source terms, or surface
parameters and checked in the same Python environment.  During emulator or surrogate
development, local RT sensitivities may help diagnose where a training set must resolve
state-space curvature or provide differentiable components for hybrid physical and
learned models.  Those uses are not demonstrated here; the present contribution is the
measured RT-kernel cost and validation needed before such workflows are built.  The
OCO-3 case below illustrates the retrieval-development use; emulator and surrogate
applications remain future work.  These uses do not require py2sess autograd VJPs to
be faster than a specialized analytic Jacobian for every production problem; they
require the sensitivities to be reproducible, checkable, and practical for the tested
development-scale prepared-input cases.

The PyTorch backend is the implemented differentiable path in this release.  A JAX
implementation could in principle provide similar VJP semantics through functional
array transformations (Bradbury et al., 2018), and XLA compilation may be attractive
for a solver written from the start as static functional array code.  A fair benchmark
would therefore require a separate JAX implementation of the BVP solve, source
integration, exceptional-row handling, and custom derivative rules, not a mechanical
translation of this PyTorch path.  The benchmark driver now includes an optional
``torch.compile'' mode so compiled PyTorch rows can be regenerated with the same
warmup, synchronization, and CSV schema as the eager rows.  The present implementation
supports this path for the thermal batch benchmark; the pseudo-spherical Solar
\texttt{include\_fo} path still relies on NumPy cached geometry precomputation and is skipped
for compiled rows.  The timings reported here remain eager-mode PyTorch unless
explicitly labeled otherwise, because compiled rows must be rerun and checked on the
target CUDA system before being mixed into the paper figures.

The OCO-3 aerosol case illustrates this development role on real data.  It required
changing aerosol vertical assumptions, phase/SSA defaults, gas-opacity caches, spectral
nuisance terms, and retrieval diagnostics while keeping the RT core fixed.  The
improvement obtained by using an external low-altitude MERRA-2 aerosol profile, rather
than a simple Gaussian pressure profile, is physically consistent with the
dust-dominated regional aerosol context.  The case also identifies the present boundary
of the demonstration.  The O$_2$ A + weak-CO$_2$ retrieval gives a plausible regional
AOD magnitude, but the strong-CO$_2$ band cannot yet be used quantitatively with the
open spectroscopy tested here, and AOD, continuum shape, vertical profile, and phase
cannot yet be inferred independently with the simple aerosol model.

The limitations are equally important.  The timing scope is prepared-input RT, not
end-to-end retrieval.  Retrieval runtime also includes spectroscopic opacity
generation, optimizer behavior, convergence criteria, prior constraints, and
instrument-specific state-vector design.  The Fortran comparison is a validated forward
benchmark; a matched Fortran sensitivity timing study would require the same prepared
RT inputs, the same repeat policy, and a clear distinction among analytic
linearization, VJP timing, and finite-difference timing.  Maximum feasible spectral-grid
size and device-memory ceilings for PyTorch VJPs remain deployment constraints that
must be assessed before applying the differentiable path to larger production
retrievals.  The synthetic sensitivity benchmark has been expanded toward a
retrieval-like prepared-input case, but any revised A100/CPU comparison should be
reported only after all matching timing assets are regenerated from the same script
revision.  The OCO-3 showcase uses
open HITRAN-based
``ABSCO-lite'' inputs rather than official OCO ABSCO tables, fixes aerosol vertical
shape and phase properties, has no strict local AERONET site collocation, and uses
MERRA-2 both as a profile-shape prior and as one contextual AOD comparison.  Additional
granules, stricter AERONET/SAM collocations, complete dust/fine aerosol optical
properties, official-quality spectroscopy, and a full retrieval error budget are needed
before making product-level aerosol claims.  Those comparisons are natural follow-up
work, but they are not needed for the present conclusion: py2sess provides a
reproducible Python pathway for 2S-ESS forward RT, CUDA execution, local
automatic-differentiation sensitivities, and constrained real-data retrieval
experiments.

\conclusions

py2sess v0.2.0 maps the published 2S-ESS RT formulation into a Python implementation
with NumPy and PyTorch backends.  It reproduces exported Fortran benchmark spectra for
large thermal-emission and solar-scattering scenes and provides reproducible
prepared-input RT timings across local CPU and CUDA systems.  Synthetic scaling shows
that forward timing is mainly controlled by wavelength count, vertical grid size, and
source-regime complexity.  PyTorch VJP sensitivities have a substantial overhead
relative to forward RT, but the overhead is nearly independent of whether one or all
local optical-property layers are active once the differentiated graph is fixed.  The
OCO-3 aerosol showcase demonstrates the same platform on measured reflected-solar
spectra, with a constrained dust-dominated O$_2$ A + weak-CO$_2$ retrieval that is
compared with MAIAC and MERRA-2 context while also exposing the spectroscopy,
continuum, and aerosol-model assumptions that prevent product-level claims.  The
result is a practical benchmarked platform for retrieval prototyping, sensitivity
studies, state-vector experimentation, and future differentiable surrogate workflows
using the 2S-ESS approximation.

\codedataavailability{The py2sess development repository is available at
\url{https://github.com/happysky19/py2sess} under the GPL-3.0-only licence.  For
submission, the exact code revision, benchmark scripts, selected figure inputs,
generated figures, exported full-spectrum input bundle, raw timing tables, and external
2S-ESS Fortran benchmark provenance should be archived in a persistent repository such
as Zenodo and cited with a DOI.  The DOI is pending in this working draft.  The local
Apple M2 Pro benchmark environment used Python 3.11.12, NumPy 2.2.6, SciPy 1.17.1,
PyTorch 2.7.1 without CUDA, and Matplotlib 3.10.3.  The CUDA Colab timing manifests
record Python 3.12.13, PyTorch 2.10.0+cu128, and CUDA 12.8 for the Tesla T4 and A100
rows.  The repository contains scripts and documentation to regenerate the
full-spectrum, synthetic forward, synthetic VJP, overhead-ratio, gradient-validation,
and OCO-3 aerosol-showcase summaries from the archived inputs.}

\authorcontribution{Author contributions will be completed before submission.}

\competinginterests{The authors declare that they have no conflict of interest.}

\begingroup
\raggedright
\begin{thebibliography}{}

\bibitem[Bell et al.(2023)]{bell2023}
Bell, E., O'Dell, C. W., Taylor, T. E., Merrelli, A., Nelson, R. R.,
Kiel, M., Eldering, A., Rosenberg, R., and Fisher, B.: Exploring bias in the
OCO-3 snapshot area mapping mode via geometry, surface, and aerosol effects,
Atmos. Meas. Tech., 16, 109--133, https://doi.org/10.5194/amt-16-109-2023,
2023.

\bibitem[Bak et al.(2021)]{bak2021}
Bak, J., Liu, X., Spurr, R., Yang, K., Nowlan, C. R., Chan Miller, C.,
Gonzalez Abad, G., and Chance, K.: Radiative transfer acceleration based on the
principal component analysis and lookup table of corrections: optimization and
application to UV ozone profile retrievals, Atmos. Meas. Tech., 14, 2659--2672,
https://doi.org/10.5194/amt-14-2659-2021, 2021.

\bibitem[Baydin et al.(2018)]{baydin2018}
Baydin, A. G., Pearlmutter, B. A., Radul, A. A., and Siskind, J. M.:
Automatic differentiation in machine learning: a survey, J. Mach. Learn. Res.,
18, 1--43, 2018.

\bibitem[Bradbury et al.(2018)]{bradbury2018}
Bradbury, J., Frostig, R., Hawkins, P., Johnson, M. J., Katariya, Y.,
Leary, C., Maclaurin, D., Necula, G., Paszke, A., VanderPlas, J.,
Wanderman-Milne, S., and Zhang, Q.: JAX: composable transformations of
Python+NumPy programs, available at: \url{https://github.com/jax-ml/jax},
2018.

\bibitem[Cao et al.(2025)]{cao2025}
Cao, L., Zhen, Z., Chen, S., and Yin, T.: A physically based differentiable
radiative transfer model (DRTM) for land surface optical and biochemical
parameters retrieval, Remote Sens. Environ., 325, 114764,
https://doi.org/10.1016/j.rse.2025.114764, 2025.

\bibitem[Chen et al.(2010)]{chen2010}
Chen, Y., Han, Y., Van Delst, P., and Weng, F.: On water vapor Jacobian in
fast radiative transfer model, J. Geophys. Res., 115, D12303,
https://doi.org/10.1029/2009JD013379, 2010.

\bibitem[Ding and Yang(2023)]{dingyang2023}
Ding, J. and Yang, P.: Tangent-linear and adjoint models for the transfer of
polarized radiation, J. Atmos. Sci., 80, 73--89,
https://doi.org/10.1175/JAS-D-22-0112.1, 2023.

\bibitem[Eldering et al.(2019)]{eldering2019}
Eldering, A., Taylor, T. E., O'Dell, C. W., and Pavlick, R.: The OCO-3
mission: measurement objectives and expected performance based on 1 year of
simulated data, Atmos. Meas. Tech., 12, 2341--2370,
https://doi.org/10.5194/amt-12-2341-2019, 2019.

\bibitem[Gelaro et al.(2017)]{gelaro2017}
Gelaro, R., McCarty, W., Suárez, M. J., Todling, R., Molod, A.,
Takacs, L., Randles, C. A., Darmenov, A., Bosilovich, M. G., Reichle, R.,
Wargan, K., Coy, L., Cullather, R., Draper, C., Akella, S., Buchard, V.,
Conaty, A., da Silva, A. M., Gu, W., Kim, G.-K., Koster, R., Lucchesi, R.,
Merkova, D., Nielsen, J. E., Partyka, G., Pawson, S., Putman, W.,
Rienecker, M., Schubert, S. D., Sienkiewicz, M., and Zhao, B.: The
Modern-Era Retrospective Analysis for Research and Applications, Version 2
(MERRA-2), J. Climate, 30, 5419--5454,
https://doi.org/10.1175/JCLI-D-16-0758.1, 2017.

\bibitem[Giles et al.(2019)]{giles2019}
Giles, D. M., Sinyuk, A., Sorokin, M. G., Schafer, J. S., Smirnov, A.,
Slutsker, I., Eck, T. F., Holben, B. N., Lewis, J. R., Campbell, J. R.,
Welton, E. J., Korkin, S. V., and Lyapustin, A. I.: Advancements in the
Aerosol Robotic Network (AERONET) Version 3 database -- automated
near-real-time quality control algorithm with improved cloud screening for Sun
photometer aerosol optical depth (AOD) measurements, Atmos. Meas. Tech., 12,
169--209, https://doi.org/10.5194/amt-12-169-2019, 2019.

\bibitem[Gordon et al.(2022)]{gordon2022}
Gordon, I. E., Rothman, L. S., Hargreaves, R. J., Hashemi, R., Karlovets,
E. V., Skinner, F. M., Conway, E. K., Hill, C., Kochanov, R. V.,
Tan, Y., Wcislo, P., Finenko, A. A., Nelson, K., Bernath, P. F.,
Birk, M., Boudon, V., Campargue, A., Chance, K. V., Coustenis, A.,
Drouin, B. J., Flaud, J.-M., Gamache, R. R., Hodges, J. T., Jacquemart, D.,
Mlawer, E. J., Nikitin, A. V., Perevalov, V. I., Rotger, M., Tennyson, J.,
Toon, G. C., Tran, H., Tyuterev, V. G., Adkins, E. M., Baker, A.,
Barbe, A., Canè, E., Császár, A. G., Dudaryonok, A., Egorov, O.,
Fleisher, A. J., Fleurbaey, H., Foltynowicz, A., Furtenbacher, T.,
Harrison, J. J., Hartmann, J.-M., Horneman, V.-M., Huang, X.,
Karman, T., Karns, J., Kassi, S., Kleiner, I., Kofman, V., Kwabia
Tchana, F., Lavrentieva, N. N., Lee, T. J., Long, D. A., Lukashevskaya,
A. A., Lyulin, O. M., Makhnev, V. Y., Matt, W., Massie, S. T.,
Melosso, M., Mikhailenko, S. N., Mondelain, D., Müller, H. S. P.,
Naumenko, O. V., Perrin, A., Polyansky, O. L., Raddaoui, E.,
Raston, P. L., Reed, Z. D., Rey, M., Richard, C., Tóbiás, R.,
Sadiek, I., Schwenke, D. W., Starikova, E., Sung, K., Tamassia, F.,
Tashkun, S. A., Vander Auwera, J., Vasilenko, I. A., Vigasin, A. A.,
Villanueva, G. L., Vispoel, B., Wagner, G., Yachmenev, A., and Yurchenko,
S. N.: The HITRAN2020 molecular spectroscopic database, J. Quant. Spectrosc.
Radiat. Transfer, 277, 107949, https://doi.org/10.1016/j.jqsrt.2021.107949,
2022.

\bibitem[Harris et al.(2020)]{harris2020}
Harris, C. R., Millman, K. J., van der Walt, S. J., Gommers, R.,
Virtanen, P., Cournapeau, D., Wieser, E., Taylor, J., Berg, S.,
Smith, N. J., Kern, R., Picus, M., Hoyer, S., van Kerkwijk, M. H.,
Brett, M., Haldane, A., del Rio, J. F., Wiebe, M., Peterson, P.,
Gerard-Marchant, P., Sheppard, K., Reddy, T., Weckesser, W.,
Abbasi, H., Gohlke, C., and Oliphant, T. E.: Array programming with NumPy,
Nature, 585, 357--362, https://doi.org/10.1038/s41586-020-2649-2, 2020.

\bibitem[Lyapustin et al.(2018)]{lyapustin2018}
Lyapustin, A., Wang, Y., Korkin, S., and Huang, D.: MODIS Collection 6 MAIAC
algorithm, Atmos. Meas. Tech., 11, 5741--5765,
https://doi.org/10.5194/amt-11-5741-2018, 2018.

\bibitem[Meador and Weaver(1980)]{meadorweaver1980}
Meador, W. E. and Weaver, W. R.: Two-stream approximations to radiative
transfer in planetary atmospheres: a unified description of existing methods and
a new improvement, J. Atmos. Sci., 37, 630--643,
https://doi.org/10.1175/1520-0469(1980)037<0630:TSATRT>2.0.CO;2, 1980.

\bibitem[Natraj et al.(2022)]{natraj2022geo}
Natraj, V., Luo, M., Blavier, J.-F., et al.: Simulated multispectral
temperature and atmospheric composition retrievals for the JPL GEO-IR Sounder,
Atmos. Meas. Tech., 15, 1251--1267, https://doi.org/10.5194/amt-15-1251-2022,
2022.

\bibitem[Natraj et al.(2023)]{natraj2023}
Natraj, V., Spurr, R., Gao, A., Le, T., Zeng, Z.-C., Fan, S., and Yung, Y. L.:
The 2 stream-exact single scattering (2S-ESS) radiative transfer model, J.
Quant. Spectrosc. Radiat. Transfer, 295, 108416,
https://doi.org/10.1016/j.jqsrt.2022.108416, 2023.

\bibitem[Natraj and Spurr(2025)]{natrajspurr2025}
Natraj, V. and Spurr, R. J. D.: Accelerating radiative transfer calculations in
the thermal infrared through principal component analysis of inherent optical
properties, J. Quant. Spectrosc. Radiat. Transfer, 339, 109415,
https://doi.org/10.1016/j.jqsrt.2025.109415, 2025.

\bibitem[Nelson and O'Dell(2019)]{nelsonodell2019}
Nelson, R. R. and O'Dell, C. W.: The impact of improved aerosol priors on
near-infrared measurements of carbon dioxide, Atmos. Meas. Tech., 12,
1495--1512, https://doi.org/10.5194/amt-12-1495-2019, 2019.

\bibitem[O'Dell et al.(2018)]{odell2018}
O'Dell, C. W., Eldering, A., Wennberg, P. O., Crisp, D., Gunson, M. R.,
Fisher, B., Frankenberg, C., Kiel, M., Lindqvist, H., Mandrake, L.,
Merrelli, A., Natraj, V., Nelson, R. R., Osterman, G. B., Payne, V. H.,
Taylor, T. E., Wunch, D., Drouin, B. J., Oyafuso, F., Chang, A.,
McDuffie, J., Smyth, M., Baker, D. F., Basu, S., Chevallier, F.,
Crowell, S. M. R., Feng, L., Palmer, P. I., Dubey, M., García, O. E.,
Griffith, D. W. T., Hase, F., Iraci, L. T., Kivi, R., Morino, I.,
Notholt, J., Ohyama, H., Petri, C., Roehl, C. M., Sha, M. K., Strong, K.,
Sussmann, R., Te, Y., Uchino, O., and Velazco, V. A.: Improved retrievals of
carbon dioxide from Orbiting Carbon Observatory-2 with the version 8 ACOS
algorithm, Atmos. Meas. Tech., 11, 6539--6576,
https://doi.org/10.5194/amt-11-6539-2018, 2018.

\bibitem[Paszke et al.(2019)]{paszke2019}
Paszke, A., Gross, S., Massa, F., Lerer, A., Bradbury, J., Chanan, G.,
Killeen, T., Lin, Z., Gimelshein, N., Antiga, L., Desmaison, A.,
Kopf, A., Yang, E., DeVito, Z., Raison, M., Tejani, A., Chilamkurthy,
S., Steiner, B., Fang, L., Bai, J., and Chintala, S.: PyTorch: an imperative
style, high-performance deep learning library, Adv. Neural Inf. Process. Syst.,
32, 8024--8035, 2019.

\bibitem[Rodgers(2000)]{rodgers2000}
Rodgers, C. D.: Inverse Methods for Atmospheric Sounding: Theory and Practice,
World Scientific, Singapore, https://doi.org/10.1142/3171, 2000.

\bibitem[Salesin et al.(2024)]{salesin2024}
Salesin, K., Knobelspiesse, K. D., Chowdhary, J., Zhai, P.-W., and Jarosz, W.:
Unifying radiative transfer models in computer graphics and remote sensing, Part
II: a differentiable, polarimetric forward model and validation, J. Quant.
Spectrosc. Radiat. Transfer, 315, 108849,
https://doi.org/10.1016/j.jqsrt.2023.108849,
2024.

\bibitem[Spurr(2006)]{spurr2006}
Spurr, R. J. D.: VLIDORT: a linearized pseudo-spherical vector discrete ordinate
radiative transfer code for forward model and retrieval studies in multilayer
multiple scattering media, J. Quant. Spectrosc. Radiat. Transfer, 102, 316--342,
https://doi.org/10.1016/j.jqsrt.2006.05.005, 2006.

\bibitem[Spurr and Natraj(2011)]{spurrnatraj2011}
Spurr, R. and Natraj, V.: A linearized two-stream radiative transfer code for
fast approximation of multiple-scatter fields, J. Quant. Spectrosc. Radiat.
Transfer, 112, 2630--2637, https://doi.org/10.1016/j.jqsrt.2011.06.014, 2011.

\bibitem[Spurr et al.(2013)]{spurr2013}
Spurr, R., Natraj, V., Lerot, C., Van Roozendael, M., and Loyola, D.:
Linearization of the principal component analysis method for radiative transfer
acceleration: application to retrieval algorithms and sensitivity studies, J.
Quant. Spectrosc. Radiat. Transfer, 125, 1--17,
https://doi.org/10.1016/j.jqsrt.2013.04.002,
2013.

\bibitem[Spurr et al.(2022)]{spurr2022}
Spurr, R., Natraj, V., Colosimo, S. F., Stutz, J., Christi, M., and Korkin, S.:
VLIDORT-QS: a quasi-spherical vector radiative transfer model, J. Quant.
Spectrosc. Radiat. Transfer, 291, 108341,
https://doi.org/10.1016/j.jqsrt.2022.108341, 2022.

\bibitem[Toon et al.(1989)]{toon1989}
Toon, O. B., McKay, C. P., Ackerman, T. P., and Santhanam, K.: Rapid
calculation of radiative heating rates and photodissociation rates in
inhomogeneous multiple scattering atmospheres, J. Geophys. Res., 94,
16287--16301, https://doi.org/10.1029/JD094iD13p16287, 1989.

\bibitem[Wang et al.(2024)]{wang2024}
Wang, Y., Kallel, A., Zhen, Z., Lauret, N., Guilleux, J., Chavanon, E., and
Gastellu-Etchegorry, J.-P.: 3D Monte Carlo differentiable radiative transfer
with DART, Remote Sens. Environ., 308, 114201,
https://doi.org/10.1016/j.rse.2024.114201, 2024.

\end{thebibliography}
\endgroup

\end{document}
